โครงงานนี้มีวัตถุประสงค์เพื่อประยุกต์ใช้กระบวนการวิทยาการข้อมูลแบบครบวงจรในการวิเคราะห์ปัจจัยก่อนการแข่งขันฟุตบอลลีกยุโรป พัฒนาแบบจำลองทำนายผลการแข่งขัน 3 ทางเลือก (เจ้าบ้านชนะ เสมอ และทีมเยือนชนะ) และตอบคำถามว่าผู้เดิมพันที่ใช้ข้อมูลสาธารณะสามารถเอาชนะตลาดเดิมพันได้หรือไม่

การดำเนินงานได้รวบรวมข้อมูล 5 ลีกใหญ่ของยุโรป ฤดูกาล 2014--2015 ถึงฤดูกาล 2026--2027 ที่กำลังแข่งขันอยู่ จาก Understat, Transfermarkt และ Football-Data.co.uk ผ่านท่อส่งข้อมูลเชิงโมดูลาร์ที่ตรวจสอบโครงสร้างข้อมูลด้วย Pandera แล้วสร้างตัวแปรก่อนแข่ง ได้แก่ คะแนน Elo ดัชนีความเหนื่อยล้า ความหนาแน่นของตารางแข่ง และฟอร์มผู้เล่น 11 ตัวจริง คัดเลือกตัวแปรแบบเดินหน้า และเปรียบเทียบแบบจำลอง Logistic Regression, Gradient-boosted Trees และแบบจำลองจำนวนประตู Poisson โดยทำนายทุกฤดูกาลจากฤดูกาลก่อนหน้าเท่านั้น จากนั้นเปรียบเทียบกับราคาของ Pinnacle ทดสอบกลยุทธ์การเดิมพัน 18 กลยุทธ์ที่ลงทะเบียนไว้ล่วงหน้าพร้อมการปรับแก้แบบ Benjamini--Hochberg และจำลองแผนการวางเงิน 5 แบบ

ผลการศึกษาพบว่า จากข้อมูลพรีเมียร์ลีกฤดูกาล 2000--2001 ถึง 2024--2025 ความได้เปรียบของเจ้าบ้าน ($+0.362$ ประตูต่อนัด) และความสัมพันธ์ของคะแนน Elo กับผลต่างประตู ($r = 0.441$) มีนัยสำคัญทางสถิติ ตัวแปรที่ช่วยทำนายได้คือคะแนน Elo และฟอร์มผู้เล่นตัวจริง ส่วนดัชนีความเหนื่อยล้าไม่ช่วยเมื่อมีคะแนน Elo อยู่แล้ว Logistic Regression ให้ผลดีที่สุด (Log Loss 0.9855, Accuracy 52.41\%) แต่ยังด้อยกว่าราคาเจ้ามือ (0.9691) การวางเดิมพันตามแบบจำลองขาดทุนประมาณร้อยละ 4--5.5 และไม่มีกลยุทธ์ใดใน 252 กรณีที่มีกำไรอย่างมีนัยสำคัญ ความได้เปรียบตามคาดหมายมีเพียงการเลือกราคาดีที่สุดระหว่างเจ้ามือ (Line Shopping) ซึ่งต้องใช้เวลาหลายฤดูกาลกว่าจะเห็นผล และเจ้ามือที่ให้ราคานั้นปิดบัญชีผู้ชนะ ส่วน Martingale และ Kelly เต็มทำให้เงินทุนหมดได้ง่าย ผลเหล่านี้ชี้ว่าตลาดเดิมพันของ 5 ลีกใหญ่มีประสิทธิภาพสูง และท่อส่งข้อมูลที่สร้างซ้ำได้นี้ใช้เป็นแนวทางประเมินกลยุทธ์การเดิมพันกีฬาอย่างเป็นระบบได้
